\section{Quantile separation and the variance lower bound}\label{sec:lower-bounds}

The remaining part of Theorem~\ref{thm:variance} is a linear lower
bound. We use a width theorem for the proper-clause model and convert
it into a second-moment estimate. Keeping the two integer quantiles
explicit makes the argument valid when the hitting-time law has atoms.

An elementary sparse-prefix estimate will also supply the positive lower
bounds on finite-size means in Appendices~\ref{sec:direct-comparison}
and~\ref{sec:appendix-c}. We record it before turning to the quantile
argument.

\begin{lemma}[A sparse satisfiable prefix]\label{lem:sparse-satisfiability}
For every real $0<c<e^{-2}$, every integer $k\ge3$, and every $n\ge k$,
put $j=\lfloor cn\rfloor$. In the independent proper-clause model,
\begin{equation}\label{eq:sparse-satisfiability}
 1-P_n(j)\le\frac{ec}{n^2}\sum_{t=1}^{\infty}t^2(e^2c)^t=O_c(n^{-2}).
\end{equation}
\end{lemma}

\begin{proof}
Form the bipartite incidence graph between the $j$ clause positions and
the $n$ variables. A matching covering all clause positions supplies a
satisfying assignment: set each matched variable to satisfy its literal
in the clause matched to it. This criterion and the deficient-set union
bound below appear in Franco and Van Gelder~\cite[Section~8]{FrancoVanGelder2003}.
By Hall's theorem~\cite[Theorem~1]{Hall1935}, failure of such a matching
implies that some set of $t$ variables contains at least $t+1$ whole
clauses. A proper $k$-clause lies in a specified $t$-set with probability
$(t)_k/(n)_k\le(t/n)^k$, interpreted as zero when $t<k$. Thus, using
$k\ge3$, $j\le cn$, and $\binom ab\le(ea/b)^b$,
\begin{align*}
 1-P_n(j)
 &\le\sum_{\substack{1\le t\le n\\t+1\le j}}
       \binom nt\binom j{t+1}(t/n)^{k(t+1)}\\
 &\le\sum_{t=1}^{n}ec(e^2c)^t(t/n)^{t+2}
 \le\frac{ec}{n^2}\sum_{t=1}^{\infty}t^2(e^2c)^t.
\end{align*}
The final series is finite because $e^2c<1$.
\end{proof}


Wilson's Corollary~4 \cite{Wilson} applies to independent uniform proper
$k$-clauses sampled with replacement. At the fixed probability levels
$3/4$ and $1/4$, its assertion for every fixed $k\ge3$ gives constants
$d_k>0$ and $n_k$ such that, for $n\ge n_k$,
\begin{equation}\label{eq:wilson-gap}
 P_n(m_1)\ge\tfrac34,\quad P_n(m_2)\le\tfrac14
 \quad\Longrightarrow\quad m_2-m_1\ge d_k\sqrt n
\end{equation}
The constants may depend on $k$. We spell out the conversion from this
quantile-width statement to a variance bound, including possible atoms.

\begin{proposition}[Linear variance lower bound]
\label{prop:variance-lower}
Let $L_n$ be positive integer caps of order $O(n)$ and suppose
$P_n(L_n-1)\le1/4$ for all sufficiently large $n$. Then
\[
 \Var(\min\{H_n,L_n\})\ge c_k n
\]
for all sufficiently large $n$, with $c_k>0$.
In particular, for every fixed $B>U_k$ and every fixed
$k\ge4$,
\begin{equation}\label{eq:sharp-variance}
 \Var(H_n)=\Theta_k(n),\qquad
 \Var(T_{n,B})=\Theta_{k,B}(n),\qquad
 \Var(\min\{H_n-1,\lfloor Bn\rfloor\})=\Theta_{k,B}(n).
\end{equation}
For $k=3$, the corresponding variances lie between positive constant
multiples of $n$ and $n\log n$.
\end{proposition}

\begin{proof}
Write $X=\min\{H_n,L_n\}$ and $G(j)=\PP(X\le j)$. Define the integer
quantiles
\[
 a=\min\{j:G(j)\ge1/4\},\qquad
 b=\min\{j:G(j)\ge3/4\}.
\]
The cap hypothesis gives $b\le L_n-1$. Thus $a-1$ and $b$ are below the
cap and, by minimality,
\[
 P_n(a-1)=1-G(a-1)>3/4,\qquad P_n(b)=1-G(b)\le1/4.
\]
Applying \eqref{eq:wilson-gap} to the endpoint probabilities above gives
$b-a\ge d_k\sqrt n-1$.
Minimality of the quantiles also gives
\[
 \PP(X\le a)\ge\tfrac14,\qquad
 \PP(X\ge b)=1-G(b-1)>\tfrac14.
\]
For an independent copy $X'$, each of the disjoint events
$\{X\le a,X'\ge b\}$ and $\{X'\le a,X\ge b\}$ has probability at
least $1/16$ once $b>a$. Therefore
\[
 \Var(X)=\tfrac12\EE(X-X')^2
 \ge\frac{(b-a)^2}{16}
 \ge\frac{d_k^2}{64}n
\]
for sufficiently large $n$.
If $L_n=\lfloor Bn\rfloor$ and $B>U_k$, then
\[
 P_n(L_n-1)\le q_k^{-2}\rho_{k,B}^n=o(1)
\]
by \eqref{eq:first-moment-tail}. The capped upper bounds proved above give the
matching upper bounds. The exact identity
\[
 \min\{H_n-1,L_n\}=\min\{H_n,L_n+1\}-1
\]
proves the lower bound for last satisfiability, since its required
condition is $P_n(L_n)\le1/4$. Finally
$\Var(H_n)\ge\Var(\min\{H_n,L_n\})$ by
Proposition~\ref{prop:cap-bridge}; Corollary~\ref{cor:uncapped} supplies
the uncapped upper bound.
This completes the variance assertions of Theorem~\ref{thm:variance}.
\end{proof}
